
\documentclass{beamer}
\usepackage[utf8]{inputenc}
\usepackage{quantikz}
\usepackage{amsmath}
\usepackage{braket}

\title{Shor's Algorithm}
\author{Morten HJ}
\date{2006}

\begin{document}

% Title Slide
\begin{frame}
    \titlepage
\end{frame}

% Outline Slide
\begin{frame}{Outline}
    \tableofcontents
\end{frame}

% Introduction Slide
\section{Introduction}
\begin{frame}{What is Shor's Algorithm?}
    \begin{itemize}
        \item Shor's algorithm is a quantum algorithm for integer factorization.
        \item It efficiently factors large integers by leveraging quantum period finding.
        \item Exponential speedup over the best-known classical algorithms.
        \item Breaks RSA encryption, which relies on the difficulty of factoring.
    \end{itemize}
\end{frame}

% Mathematical Background
\section{Mathematical Background}
\begin{frame}{Number Theory Basics}
    \begin{itemize}
        \item The problem: Factor an integer $N$ into its prime factors.
        \item Reduce factoring to period finding:  
        \[ f(x) = a^x \mod N, \]
        where $a$ is randomly chosen such that $\gcd(a, N) = 1$.
        \item The period $r$ satisfies:
        \[ a^r \equiv 1 \mod N. \]
        \item Once $r$ is known, factors are given by:
        \[ \gcd(a^{r/2} \pm 1, N). \]
    \end{itemize}
\end{frame}

% Quantum Period Finding
\section{Quantum Period Finding}
\begin{frame}{Quantum Period Finding: Key to Shor's Algorithm}
    \begin{itemize}
        \item Quantum subroutine for finding the period $r$ of $f(x)$.
        \item Similar to Quantum Phase Estimation (QPE).
        \item Uses two quantum registers:
        \begin{itemize}
            \item First register: Superposition of all possible inputs $x$.
            \item Second register: Computes $f(x) = a^x \mod N$.
        \end{itemize}
    \end{itemize}
\end{frame}

% Quantum Circuit Design
\section{Quantum Circuit Design}
\begin{frame}{Quantum Circuit for Shor's Algorithm}
    \begin{figure}
        \centering
        \begin{quantikz}
            \lstick{$\ket{0}^{\otimes n}$} & \qwbundle{n} & \gate{H} & \ctrl{1} & \gate{QFT^\dagger} & \meter{} \\
            \lstick{$\ket{1}$} & \qw & \qw & \gate{U_f} & \qw & \qw
        \end{quantikz}
        \caption{Quantum circuit for period finding in Shor's algorithm.}
    \end{figure}
\end{frame}

\begin{frame}{Applying Quantum Fourier Transform (QFT)}
    \begin{itemize}
        \item The QFT is applied to the first register after the controlled unitary operations.
        \item Peaks in the measurement correspond to multiples of $1/r$.
        \item Post-processing classically determines $r$ by continued fractions.
    \end{itemize}
\end{frame}

% Classical Post-Processing
\section{Classical Post-Processing}
\begin{frame}{Classical Post-Processing}
    \begin{itemize}
        \item Once $r$ is found, factors of $N$ are computed by:
        \[ \gcd(a^{r/2} \pm 1, N). \]
        \item If $r$ is odd or $a^{r/2} \equiv -1 \mod N$, try a different $a$.
        \item With high probability, this yields a non-trivial factor of $N$.
    \end{itemize}
\end{frame}

% Complexity and Practical Considerations
\section{Complexity and Practical Considerations}
\begin{frame}{Complexity and Practical Considerations}
    \begin{itemize}
        \item Quantum Complexity: Polynomial in $\log N$.
        \item Classical Best Known: Sub-exponential (Number Field Sieve).
        \item Quantum Advantage: Exponential speedup over classical algorithms.
        \item Practical Challenges: Quantum error correction, large qubit counts.
    \end{itemize}
\end{frame}

% Conclusion
\section{Conclusion}
\begin{frame}{Conclusion}
    \begin{itemize}
        \item Shor's algorithm is a groundbreaking quantum algorithm for factoring.
        \item Combines quantum period finding with classical number theory.
        \item Highlights the potential of quantum computing to break RSA.
    \end{itemize}
\end{frame}

% References
\begin{frame}{References}
    \begin{itemize}
        \item P. Shor, "Algorithms for Quantum Computation: Discrete Logarithms and Factoring".
        \item M. A. Nielsen and I. L. Chuang, \textit{Quantum Computation and Quantum Information}.
    \end{itemize}
\end{frame}

\end{document}
